\section{Predecessor and Successor in a BST}
\label{sec:bst-pred-succ}

\problemstatement{Given the root of a BST and a key $X$ (which may or may
not itself be present in the tree), find the \textbf{inorder predecessor}
of $X$ (the largest value in the tree strictly less than $X$) and the
\textbf{inorder successor} of $X$ (the smallest value strictly greater
than $X$).}

\begin{patternbox}
Predecessor and successor are floor and ceil
(Section~\ref{sec:bst-floor-ceil}) with \textbf{strict} inequalities
instead of non-strict ones -- the keyword difference between ``floor''
(largest $\le X$) and ``predecessor'' (largest $<X$) is exactly the
difference between $\ge$ and $>$ that distinguished lower bound from
upper bound back in the Binary Search chapter.
\end{patternbox}

\subsection*{Understanding the problem carefully}

The walk is identical in shape to floor/ceil: at each node, compare to
$X$ and decide which candidate to update and which direction to move.
The only change is that when the current node's value equals $X$
\emph{exactly}, we must \textbf{not} stop immediately (unlike floor/ceil,
where an exact match ends the search for both) -- a predecessor or
successor is required to be \emph{strictly} different from $X$, so if the
current node's value equals $X$, we must continue descending: moving left
to find a predecessor strictly smaller, or right to find a successor
strictly larger.

\begin{ideabox}[Same Walk as Floor/Ceil, but Treat an Exact Match as ``Keep Going'']
For the predecessor: at each node, if its value $<X$ (this includes the
case of looking for a value strictly less, so equality to $X$ does
\emph{not} qualify): record it as the current best predecessor candidate,
and move right (search for an even closer one). If its value $\ge X$: do
not record anything; move left. Symmetrically for the successor, with
roles of $<$/$\ge$ swapped to $>$/$\le$.
\end{ideabox}

\subsection*{Why treating equality as ``move past it'' (not ``stop'') is the only change needed}

In floor/ceil, an exact match is itself simultaneously the floor and the
ceil, so the search correctly terminates there. In predecessor/successor,
an exact match is explicitly excluded from being its own answer, so the
search must instead behave exactly as if that node were just slightly
too small (for predecessor purposes, treat it like any other ``$\ge X$''
node and move left without recording it) or slightly too large (for
successor purposes, move right without recording it) -- everything else
about the boundary-narrowing logic carries over unchanged from
Section~\ref{sec:bst-floor-ceil}.

\subsection*{Algorithm}

\begin{enumerate}
  \item \textsc{Predecessor}: $\textit{pred} \gets$ undefined,
        $\textit{node} \gets \textit{root}$. While $\textit{node}$ is
        not null: if $\textit{node}.\textit{val} < X$:
        $\textit{pred} \gets \textit{node}.\textit{val}$; move right.
        Else: move left.
  \item \textsc{Successor}: symmetric, with $>$ in place of $<$ and
        left/right swapped.
\end{enumerate}

\subsection*{Dry run}

\dryrun{Take the BST with root $8$, left child $3$ (children $1,6$, where
$6$ has left child $4$), right child $10$ (children $9,14$). Query
$X=6$.}

\textbf{Predecessor:} At $8$: $8<6$? No, move left to $3$. At $3$:
$3<6$? Yes, $\textit{pred}=3$, move right to $6$. At $6$: $6<6$? No
(equal, not strictly less), move left to $4$. At $4$: $4<6$? Yes,
$\textit{pred}=4$, move right (null). Stop. Predecessor $=4$.

\textbf{Successor:} At $8$: $8>6$? Yes, $\textit{succ}=8$, move left to
$3$. At $3$: $3>6$? No, move right to $6$. At $6$: $6>6$? No (equal),
move right (null). Stop. Successor $=8$.

\subsection*{C++ implementation}

\begin{lstlisting}
#include <bits/stdc++.h>
using namespace std;

struct TreeNode {
    int val;
    TreeNode* left;
    TreeNode* right;
    TreeNode(int v) : val(v), left(nullptr), right(nullptr) {}
};

int findPredecessor(TreeNode* root, int X) {
    int pred = -1;
    while (root) {
        if (root->val < X) {
            pred = root->val;
            root = root->right;
        } else {
            root = root->left;
        }
    }
    return pred;
}

int findSuccessor(TreeNode* root, int X) {
    int succ = -1;
    while (root) {
        if (root->val > X) {
            succ = root->val;
            root = root->left;
        } else {
            root = root->right;
        }
    }
    return succ;
}

int main() {
    TreeNode* root = new TreeNode(8);
    root->left = new TreeNode(3);
    root->left->left = new TreeNode(1);
    root->left->right = new TreeNode(6);
    root->left->right->left = new TreeNode(4);
    root->right = new TreeNode(10);
    root->right->left = new TreeNode(9);
    root->right->right = new TreeNode(14);

    cout << "Predecessor: " << findPredecessor(root, 6) << endl;
    cout << "Successor: " << findSuccessor(root, 6) << endl;
    return 0;
}
\end{lstlisting}

\begin{complexitybox}
\textbf{Time Complexity.} $O(h)$ for each of predecessor and successor.
\textbf{Space Complexity.} $O(1)$ for the iterative version shown.
\end{complexitybox}

\begin{takeawaybox}
Predecessor/successor and floor/ceil are the same boundary-search
template; the only behavioral difference is what happens on an exact
match (stop immediately for floor/ceil, keep narrowing for
predecessor/successor). When given a problem statement, check carefully
whether equality should count as a valid answer or must be strictly
excluded -- that single detail is the entire distinction between these
two near-identical problems.
\end{takeawaybox}
